% =========================================================
% Satisfaction, Truth, Validity, Satisfiability
% =========================================================

\section{Semantics: Satisfaction and Truth}

\begin{remark*}[Section toolkit: Satisfaction ledger]
Satisfaction is the recursive semantic relation for formulas. Formulas with
free variables are evaluated relative to an assignment. Sentences have no free
variables, so their truth in a structure does not depend on which assignment is
chosen.
\end{remark*}

\begin{definitionbox}{Definition (Satisfaction)}
\begin{definition}[Satisfaction]
\label{def:fol-satisfaction}
Let \(\mathcal M=\langle D,I\rangle\) be a structure, let \(s\) be a variable
assignment, and let \(\varphi\) be a well-formed formula. The relation
\(\mathcal M,s\models\varphi\) is defined recursively by the usual clauses:
atomic formulas are evaluated by term values and interpreted predicate
relations; equality compares term values; connectives follow the truth
conditions for \(\neg,\land,\lor,\to,\leftrightarrow\); and
\[
\mathcal M,s\models\forall x\,\psi
\Longleftrightarrow
\forall d\in D\;(\mathcal M,s[x\mapsto d]\models\psi),
\]
\[
\mathcal M,s\models\exists x\,\psi
\Longleftrightarrow
\exists d\in D\;(\mathcal M,s[x\mapsto d]\models\psi).
\]
\end{definition}
\end{definitionbox}

\begin{remark*}[Standard quantified statement]
\[
\begin{aligned}
&\mathcal M,s\models P(t_1,\ldots,t_n)
\Longleftrightarrow
(\llbracket t_1\rrbracket_{\mathcal M,s},\ldots,
\llbracket t_n\rrbracket_{\mathcal M,s})\in I(P),\\
&\mathcal M,s\models t_1=t_2
\Longleftrightarrow
\llbracket t_1\rrbracket_{\mathcal M,s}
=
\llbracket t_2\rrbracket_{\mathcal M,s},\\
&\mathcal M,s\models\neg\psi
\Longleftrightarrow
\mathcal M,s\not\models\psi,\\
&\mathcal M,s\models(\psi\land\chi)
\Longleftrightarrow
\mathcal M,s\models\psi \text{ and } \mathcal M,s\models\chi,\\
&\mathcal M,s\models\forall x\,\psi
\Longleftrightarrow
\forall d\in D\;(\mathcal M,s[x\mapsto d]\models\psi),\\
&\mathcal M,s\models\exists x\,\psi
\Longleftrightarrow
\exists d\in D\;(\mathcal M,s[x\mapsto d]\models\psi).
\end{aligned}
\]
\end{remark*}

\begin{remark*}[Definition predicate reading]
\(\operatorname{Satisfaction}(M,s,\varphi)\) means
\(\mathcal M,s\models\varphi\).
\end{remark*}

\begin{remark*}[Negated quantified statement]
\[
\mathcal M,s\not\models\varphi.
\]
\end{remark*}

\begin{remark*}[Negation predicate reading]
\(\neg\operatorname{Satisfaction}(M,s,\varphi)\) means that \(\varphi\) is not
satisfied in \(M\) under \(s\).
\end{remark*}

\begin{remark*}[Failure modes]
Satisfaction fails according to the outermost formation rule of the formula:
an atomic tuple can miss the interpreted relation, equality can compare
different term values, a connective can have the wrong truth pattern, or a
quantifier can have a counterexample or lack a witness.
\end{remark*}

\begin{remark*}[Failure mode decomposition]
\[
\underbrace{(\llbracket t_i\rrbracket)\notin I(P)}_{\text{atomic failure}}
\quad
\underbrace{\llbracket t_1\rrbracket\neq\llbracket t_2\rrbracket}_{\text{equality failure}}
\quad
\underbrace{\exists d\in D\;(\mathcal M,s[x\mapsto d]\not\models\psi)}_{\text{universal failure}}.
\]
\end{remark*}

\begin{remark*}[Interpretation]
Satisfaction is the bridge from grammar to semantic use. It is recursive
because formulas are recursive. The assignment is essential exactly where free
variables can affect the value of the formula.
\end{remark*}

\begin{dependencies}
\begin{itemize}
  \item \hyperref[def:structure]{Structure}
  \item \hyperref[def:var-assign]{Variable Assignment}
  \item \hyperref[def:mod-assign]{Modified Assignment}
  \item \hyperref[def:term-interp]{Term Evaluation}
  \item \hyperref[def:wff-pred]{Well-Formed Formula}
\end{itemize}
\end{dependencies}

\begin{definitionbox}{Definition (Truth in a Structure)}
\begin{definition}[Truth in a Structure]
\label{def:truth-structure}
Let \(\varphi\) be a sentence and let \(\mathcal M\) be a structure. The
sentence \(\varphi\) is \emph{true in} \(\mathcal M\), written
\(\mathcal M\models\varphi\), exactly when
\(\mathcal M,s\models\varphi\) for every variable assignment \(s\) for
\(\mathcal M\).
\end{definition}
\end{definitionbox}

\begin{remark*}[Standard quantified statement]
\[
\mathcal M\models\varphi
\Longleftrightarrow
\forall s:\mathsf{Var}_{\mathcal L}\to D\;(\mathcal M,s\models\varphi),
\qquad
\varphi \text{ a sentence.}
\]
\end{remark*}

\begin{remark*}[Definition predicate reading]
\(\operatorname{TruthInStructure}(M,\varphi)\) means
\(\mathcal M\models\varphi\).
\end{remark*}

\begin{remark*}[Negated quantified statement]
\[
\mathcal M\not\models\varphi
\Longleftrightarrow
\exists s:\mathsf{Var}_{\mathcal L}\to D\;(\mathcal M,s\not\models\varphi).
\]
\end{remark*}

\begin{remark*}[Negation predicate reading]
\(\neg\operatorname{TruthInStructure}(M,\varphi)\) means that \(\varphi\) is
false in \(M\).
\end{remark*}

\begin{remark*}[Failure modes]
Truth in a structure fails when some assignment makes the sentence fail. For
sentences, this is independent of assignment, but the quantified form keeps the
definition uniform with satisfaction.
\end{remark*}

\begin{remark*}[Failure mode decomposition]
\[
\underbrace{\exists s}_{\text{assignment witness}}
\quad
\underbrace{\mathcal M,s\not\models\varphi}_{\text{satisfaction failure}}.
\]
\end{remark*}

\begin{remark*}[Interpretation]
Sentences have no free variables. Therefore the particular assignment does not
carry mathematical information once the formula is a sentence; the notation
\(\mathcal M\models\varphi\) suppresses it.
\end{remark*}

\begin{dependencies}
\begin{itemize}
  \item \hyperref[def:fol-satisfaction]{Satisfaction}
\end{itemize}
\end{dependencies}

\begin{definitionbox}{Definition (Validity)}
\begin{definition}[Validity]
\label{def:validity}
A formula \(\varphi\) is \emph{valid}, or \emph{logically true}, exactly when
\(\mathcal M,s\models\varphi\) for every structure \(\mathcal M\) for the
language of \(\varphi\) and every variable assignment \(s\) for \(\mathcal M\).
\end{definition}
\end{definitionbox}

\begin{remark*}[Standard quantified statement]
\[
\models\varphi
\Longleftrightarrow
\forall \mathcal M\;\forall s\;(\mathcal M,s\models\varphi).
\]
\end{remark*}

\begin{remark*}[Definition predicate reading]
\(\operatorname{ValidFormula}(\varphi)\) means that \(\varphi\) is satisfied in
every structure under every assignment.
\end{remark*}

\begin{remark*}[Negated quantified statement]
\[
\not\models\varphi
\Longleftrightarrow
\exists \mathcal M\;\exists s\;(\mathcal M,s\not\models\varphi).
\]
\end{remark*}

\begin{remark*}[Negation predicate reading]
\(\neg\operatorname{ValidFormula}(\varphi)\) means that some structure and
assignment make \(\varphi\) fail.
\end{remark*}

\begin{remark*}[Failure modes]
Validity fails by a counterexample structure together with an assignment.
\end{remark*}

\begin{remark*}[Failure mode decomposition]
\[
\underbrace{\mathcal M}_{\text{counterexample structure}}
\qquad
\underbrace{s}_{\text{counterexample assignment}}
\qquad
\underbrace{\mathcal M,s\not\models\varphi}_{\text{failure of satisfaction}}.
\]
\end{remark*}

\begin{remark*}[Interpretation]
Validity is truth by logic alone. It does not depend on a special structure,
special domain, or special interpretation of the non-logical symbols.
\end{remark*}

\begin{dependencies}
\begin{itemize}
  \item \hyperref[def:fol-satisfaction]{Satisfaction}
\end{itemize}
\end{dependencies}

\begin{definitionbox}{Definition (Satisfiability)}
\begin{definition}[Satisfiability]
\label{def:fol-satisfiability}
A formula \(\varphi\) is \emph{satisfiable} exactly when there are a structure
\(\mathcal M\) and a variable assignment \(s\) such that
\(\mathcal M,s\models\varphi\).
\end{definition}
\end{definitionbox}

\begin{remark*}[Standard quantified statement]
\[
\varphi \text{ is satisfiable}
\Longleftrightarrow
\exists \mathcal M\;\exists s\;(\mathcal M,s\models\varphi).
\]
\end{remark*}

\begin{remark*}[Definition predicate reading]
\(\operatorname{SatisfiableFormula}(\varphi)\) means that \(\varphi\) is
satisfied in at least one structure under at least one assignment.
\end{remark*}

\begin{remark*}[Negated quantified statement]
\[
\varphi \text{ is unsatisfiable}
\Longleftrightarrow
\forall \mathcal M\;\forall s\;(\mathcal M,s\not\models\varphi).
\]
\end{remark*}

\begin{remark*}[Negation predicate reading]
\(\neg\operatorname{SatisfiableFormula}(\varphi)\) means that no structure and
assignment satisfy \(\varphi\).
\end{remark*}

\begin{remark*}[Failure modes]
Satisfiability fails only globally: every attempted structure and assignment
fails to satisfy the formula.
\end{remark*}

\begin{remark*}[Failure mode decomposition]
\[
\underbrace{\forall \mathcal M}_{\text{all structures}}
\quad
\underbrace{\forall s}_{\text{all assignments}}
\quad
\underbrace{\mathcal M,s\not\models\varphi}_{\text{no satisfying instance}}.
\]
\end{remark*}

\begin{remark*}[Interpretation]
Satisfiability asks whether a formula can be made true somewhere. This is the
semantic counterpart of consistency only in the minimal Volume I sense; deeper
proof-theoretic distinctions are deferred.
\end{remark*}

\begin{dependencies}
\begin{itemize}
  \item \hyperref[def:fol-satisfaction]{Satisfaction}
\end{itemize}
\end{dependencies}
